\section{Divisor-defect feedback from the actual residual (NS-75)}
\label{sec:ns75-divisor-feedback}

\paragraph{Question.}
The fixed M\"obius templates in NS-73 do not use the optimized old
coefficients. We now test a correction that does: cancel the next block
of derivative jumps of the actual residual. This yields an exact
arithmetic construction, but removing those local jumps is not a
norm-descent theorem.

Write the complete residual as
$R_N=\tau-\sum_{n\leq N}c_n\sigma_n$, with the actual optimized $c_n$,
and put
\[
 \eta_c=\sum_{n\leq N}\frac{c_n}{n},\qquad
 r_N(m)=\mathbf1_{m=1}+\sum_{\substack{n\mid m\\n\leq N}}c_n.
\]
The following statements also hold for arbitrary fixed real old coefficients.

\begin{proposition}[Complete divisor-log cells and derivative jumps]
\label{prop:ns75-divisor-cells}
For every $x>0$,
\begin{equation}
 R_N(x)=-\frac{\eta_c}{x}
       +\sum_{m\leq1/x}r_N(m)\log\frac1{mx}.
 \label{eq:ns75-divisor-cells}
\end{equation}
At $x=1/m$ the function is continuous and the right-minus-left jump of
its first derivative is $m r_N(m)$. In distributions on $(0,\infty)$,
\begin{equation}
 R_N''=\frac1{x^2}\sum_{m\leq1/x}r_N(m)
           -\frac{2\eta_c}{x^3}
           +\sum_{m\geq1}m r_N(m)\delta_{1/m}.
 \label{eq:ns75-jump-measure}
\end{equation}
The atomic sum is locally finite away from zero. The complete exterior
tail is $R_N(x)=-\eta_c/x$ for $x>1$.
\end{proposition}
\begin{proof}
The exact identity
$A(z)=z-\sum_{k\leq z}\log(z/k)$ gives
\[
 \sum_{n\leq N}c_n\sigma_n(x)
 =\frac{\eta_c}{x}-\sum_{m\leq1/x}
       \left(\sum_{\substack{n\mid m\\n\leq N}}c_n\right)
       \log\frac1{mx}.
\]
For each fixed $x$ both sums are finite, so this rearrangement requires
no limiting interchange. The target is exactly the term $m=1$ in the
logarithmic sum. A term entering at $1/m$ has value zero there, proving
continuity. Between endpoints, differentiation gives
$R_N'=\eta_c/x^2-x^{-1}\sum_{m\leq1/x}r_N(m)$.
The stated jump and second-derivative distribution follow, including
the target endpoint $m=1$.
\end{proof}

\paragraph{An exact local cancellation step.}
For $N<n\leq2N$ choose physical new coefficients
\[
 a_n^{(p)}=-r_N(n)\log^p(2N/n),\qquad p=0,1,2,
\]
with the constant taper convention for $p=0$.
With $p=0$, append these coefficients while holding all old coefficients
fixed. For $N<m\leq2N$, every proper divisor of $m$ is at most $N$.
Consequently the updated divisor defect is exactly
$r_N(m)+a_m^{(0)}=0$ throughout this next block. The old defects are
unchanged. Other cells, the regular part of~\eqref{eq:ns75-jump-measure}
and both tails remain in the norm calculation. No smallness follows
from cancellation of the displayed atomic part alone.

If the old coefficients happen to be $c_n=-\mu(n)$, then
$r_N(m)=\mu(m)$ on this next block, so the constant rule specializes
to NS-73's sharp M\"obius template. This is a specialization, not an
assumption about the optimized old coefficients.

For comparison with the optimum we also convert each $a^{(p)}$ to
difference coordinates using~\eqref{eq:ns73-coordinate-change}, apply
the complete old-space complement, and optimize its scalar step.
The exact gain and the best three-direction span gain are then those
of Proposition~\ref{prop:ns73-span}. This old-space reoptimization
can change the local jumps again; we do not claim it preserves the
raw cancellation property. We separately record the complete error
of the held-old, unit-sized step:
\[
 \frac{\|R_N-\sum_{n=N+1}^{2N}a_n^{(p)}\sigma_n\|_2^2}{E_N}
 =1+\frac{(a^{(p)})^TG_{\rm new,new}a^{(p)}
                 -2(g^Tv^{(p)})}{E_N}.
\]
Old-space orthogonality justifies the numerator equality here;
the raw physical Gram in this expression is not replaced by a
projected or model Gram.

\paragraph{Certified finite outcome.}
The complete 256/384-bit replays and doubled-cutoff replay at $N=256$
give the following strict enclosures for the fraction of full block gain.
The single cancellation direction is optimally scaled \emph{after} old-space
projection; its scalar is not constrained positive.
\begin{center}
\begin{tabular}{rccc}
\hline
$N$&Cancellation direction&Best three-direction span&Diagonal curvature\\
\hline
16&(.3110,.3111)&(.5836,.5837)&(.9308,.9310)\\
32&(.4292,.4294)&(.8086,.8087)&(.9040,.9041)\\
64&(.6954,.6956)&(.9063,.9065)&(.9783,.9784)\\
128&(.0029,.0030)&(.7232,.7233)&(.9852,.9854)\\
256&(.0010,.0011)&(.5602,.5603)&(.9763,.9764)\\
\hline
\end{tabular}
\end{center}
At $N=256$ the span retains $56.02048\%$ of full gain, versus
$97.63068\%$ for curvature. The raw held-old cancellation step increases
the complete squared error by a factor in $(1495.54,1495.55)$.
Indeed all fifteen held-old unit steps tested here increase squared error.
The cancellation direction's correlation at $N=256$ is negative, so its
best scalar step reverses the prescribed sign. Merely cancelling derivative
jumps is therefore not a descent principle even on these certified blocks.
This is a finite failure of the proposed unit step and finite underperformance
of its span, not an asymptotic exclusion of every divisor-feedback scheme.
Eleven exact Maxima checks cover the identities and coordinate conversion;
all finite gains retain complete tails, the true old projection and all
three-direction cross terms.

\section{The continuous relaxation and its exact inner-factor defect (NS-76)}
\label{sec:ns76-inner-distance}

\paragraph{Scope and prior work.}
Burnol's quantitative Nyman formula identifies the $q=1$ continuous
approximation defect through the off-critical Blaschke product
\cite[Theorem 1.2 and Lemma 2.4]{BurnolNyman2001}. We apply the same
classical Hardy-space mechanism to this project's $q=2$ target, keeping
its full exterior tail. The calculation gives an exact distance for
\emph{continuous} dilations and a lower bound for integer dilations.
No unconditional equality of the two nonzero distances is assumed.
This is an explicit reformulation of the unknown arithmetic inner
factor, not a new zero-free estimate.

Let $\mathcal T$ be the complete tail space from NS-74, and let
$\mathbb H^2$ denote the Hardy space of $\Re s>1/2$, with boundary norm
$(2\pi)^{-1}\int_{\mathbb R}|f(1/2+it)|^2dt$.
Multiplication on the Mellin boundary by $(s-1)/s$ induces a unitary
map $W$ on the full $L^2$ space. It maps $\mathcal T$ onto
$L^2(0,1)$, whose Mellin image is $\mathbb H^2$.
Indeed in logarithmic coordinates it is
\[
 (Wu)(t)=u(t)-\int_{-\infty}^t e^{-(t-r)/2}u(r)\,dr.
\]
For $t<0$ a tail vector is $C e^{t/2}$ and cancels exactly in this
formula. Conversely the adjoint applied to a function supported on
$t\geq0$ has on $t<0$ a constant multiple of $e^{t/2}$.
Thus the map is onto, and the full tail has been accounted for.

The transformed target and first atom have Mellin transforms
\[
 T(s)=\frac{s-1}{s^3},\qquad
 F(s)=-\frac{(s-1)\zeta(s)}{s^3}.
\]
Write $B$ for the off-critical Blaschke product, normalized by $B(1)>0$:
\[
 B(s)=\prod_{\substack{\zeta(\rho)=0\\\Re\rho>1/2}}
 \frac{s-\rho}{s-(1-\overline\rho)}
 \frac{1-\overline\rho}{\rho}
 \left|\frac{\rho}{1-\rho}\right|.
\]
Multiplicities are included. Burnol's factorization proves that the
inner factor of $(s-1)\zeta(s)/s^2$ is exactly $B$: analytic
continuation excludes a finite-boundary singular factor, and the
real-axis asymptotic excludes an exponential factor. The additional
factor $1/s$ is outer, so $F$ has the same inner factor. The product
and its derivatives converge locally in $\Re s>1/2$ away from its
zeros; $B(1)$ and its first two derivatives are real by conjugation
symmetry.

\begin{proposition}[Exact distance for the complete continuous family]
\label{prop:ns76-continuous-distance}
Let $\mathcal V_{\rm cont}$ be the closed span in the original Hilbert
space of $\sigma_u$, for all real $u\geq1$. Its image under $W$ and
Mellin transformation is $B\mathbb H^2$.
Put $b=B(1)$, $d=B'(1)$ and $e=B''(1)$. Then
\begin{equation}
 D_{\rm cont}^2:=\operatorname{dist}(\tau,\mathcal V_{\rm cont})^2
 =2-2b^2+2bd-d^2-\frac{e^2}{4}.
 \label{eq:ns76-continuous-distance}
\end{equation}
This expression is nonnegative and at most every integer finite
squared error $E_N$. It vanishes exactly when RH holds.
The formula contains unknown off-critical zero information and
provides no unconditional upper bound tending to zero.
\end{proposition}
\begin{proof}
The normalized real dilations of $\sigma_1$ become the causal
translation semigroup applied to $F$. The Beurling--Lax theorem gives
that its closed generated subspace is the inner factor $B$ times
$\mathbb H^2$. This use of the full continuous semigroup is the reason
we do not replace it by the integer set without a separate argument.

For $\Re z>1/2$, let $k_z(s)=1/(s+\overline z-1)$ be the evaluation
kernel. Multiplication $M_B$ is an isometry and
$M_B^*k_z=\overline{B(z)}k_z$.
Differentiating at real $z=1$ is legitimate in the Hilbert norm,
as is seen directly from the integrable kernel derivatives. Since
$T=-\partial_z k_z|_1-\tfrac12\partial_z^2 k_z|_1$, it follows that
\[
 M_B^*T=-\frac{d+e/2}{s}+\frac{b+d}{s^2}-\frac{b}{s^3}.
\]
The complete boundary Gram for $(1/s,1/s^2,1/s^3)$ is
\[
 \begin{pmatrix}1&1&1\\1&2&3\\1&3&6\end{pmatrix};
\]
its entries follow by integrating the Laplace representatives
$e^{-t/2}t^{j-1}/(j-1)!$ on $t>0$.
Therefore $\|M_B^*T\|^2=2b^2-2bd+d^2+e^2/4$.
The orthogonal projection is $B M_B^*T$, with that same norm, and
$\|T\|^2=\|\tau\|^2=2$. This proves the distance identity and
nonnegativity without any sign assumption on $d$ or $e$.

The integer span is contained in the continuous span, giving
$E_N\geq D_{\rm cont}^2$. If RH holds, the product is empty and
$(b,d,e)=(1,0,0)$. Conversely, any off-critical zero $\rho$ gives a
zero of $B$ in this half-plane, whereas $T(\rho)\ne0$ since
$\rho\notin\{0,1\}$. Continuous evaluation in $\mathbb H^2$ then
prevents $T$ from belonging to the closed subspace $B\mathbb H^2$,
so its distance is positive.
\end{proof}

\begin{corollary}[The fixed component invisible to every block gain]
\label{cor:ns76-invisible-component}
Let $P_{\rm cont}$ and $P_N$ be the complete projections onto the
continuous and finite integer spaces. Then exactly
\[
 R_N=(I-P_{\rm cont})\tau+(P_{\rm cont}\tau-P_N\tau),
 \qquad
 E_N=D_{\rm cont}^2+\|P_{\rm cont}\tau-P_N\tau\|^2.
\]
Every new correlation $g_n$ and every selected-direction gain depend
only on the second component. The first is orthogonal to every real
dilation, hence also to each projected new integer atom.
\end{corollary}
\begin{proof}
The finite integer space is a subspace of $\mathcal V_{\rm cont}$,
so $P_NP_{\rm cont}=P_N$. The two displayed components are
orthogonal. Taking inner products with any old-projected new atom
annihilates the first, proving the final statement.
\end{proof}

\begin{proposition}[The integer and continuous closures are strictly different]
\label{prop:ns76-strict-closure}
Let $\mathcal V_{\rm int}$ be the closed span of the integer atoms.
Then $\mathcal V_{\rm int}\subsetneq\mathcal V_{\rm cont}$,
unconditionally. More explicitly,
\[
 \operatorname{dist}(\sigma_{3/2},\mathcal V_{\rm int})^2>.00111713.
\]
This is a bound for $\sigma_{3/2}$, not for the target $\tau$.
\end{proposition}
\begin{proof}
On $I=(1/2,1)$, the first integer atom is
$\sigma_1=1/x+\log x$ and every integer atom with $n\geq2$
is $1/(nx)$. Their restricted closed span lies in the closed
two-dimensional space $\operatorname{span}\{1/x,\log x\}$.
In contrast, $\sigma_u$ for $1<u<2$ has a nonzero derivative jump
at the interior point $x=1/u$; it cannot equal an element of that
analytic two-dimensional span, even almost everywhere. This already
proves strict inclusion.

For the quantitative bound, put $u=3/2$, $L=\log2$, $c=\log u$ and
$h=L-c$. The complete local Gram, loads and squared norm are
\[
 G_I=\begin{pmatrix}1&-L^2/2\\-L^2/2&1-L-L^2/2\end{pmatrix},\qquad
 b_I=\begin{pmatrix}
 u^{-1}-h^2/2\\
 (c+2-L^2/2)/u+c(L+1)/2-(L^2+2L+2)/2
 \end{pmatrix},
\]
\[
 q_I=u^{-2}-h^2/u+2/u-h^2/2-h-1.
\]
These follow by splitting only at $x=1/u$, with
$\sigma_u=1/(ux)+\log(ux)$ on the left and $1/(ux)$ on the right.
The Gram is positive definite. Its exact local projection error
$q_I-b_I^TG_I^{-1}b_I$ lies in $(.00111713,.00111714)$ by independent
256/384-bit enclosures; Maxima checks all complete elementary integrals.
Restriction of a norm to $I$ can only decrease it, so this local
minimum bounds the full-space distance below. No omitted tail is
set to zero in the asserted full-space inequality.
\end{proof}

\paragraph{What the calculation does and does not resolve.}
The $q=2$ continuous defect is now expressed through three values of the
classical unknown inner factor. It is not estimated away. In particular,
the integer limiting error may contain an additional approximation
component. The closures are strictly different, but that fact alone
does not decide whether their two distances from this particular target
are different. No equality of positive target distances is supplied.
The original sufficient cofinal contraction remains valid, because it
would force the entire $E_N$ to zero, including this fixed component.
Replacing $E_N$ by $E_N-D_{\rm cont}^2$ would remove exactly the part
that must be ruled out and would not prove RH.

\section{How a high zero tail can hide the continuous defect (NS-77)}
\label{sec:ns77-continuous-tail}

\paragraph{Scope.}
We give a complete upper bound on the continuous defect in NS-76 using
positive sums over hypothetical off-critical zeros. This quantifies a
limitation of finite geometric evidence. It is not an upper bound on
$E_N$ or its integer limit, not a new verified zero height, and not a
proof that the defect is exactly zero.

Suppose every off-critical zero in the upper half-plane has height
$\gamma>H\geq1$. Index the conjugate pairs by $\rho=\beta+i\gamma$,
with $1/2<\beta<1$, once per pair and with multiplicity, and put
\[
 S_2=\sum_\rho\frac{2\beta-1}{\gamma^2},\qquad
 S_4=\sum_\rho\frac{2\beta-1}{\gamma^4}.
\]
These sums converge by the usual zero-counting bound. Empty sums are zero.

\begin{proposition}[A complete positive zero-tail budget]
\label{prop:ns77-tail-budget}
For the continuous $q=2$ distance,
\begin{equation}
 0\leq D_{\rm cont}^2\leq
       \min\left\{2,\frac83S_2^3+5S_4\right\}.
 \label{eq:ns77-tail-budget}
\end{equation}
If the total upper-half-plane zero count satisfies
$N(t)\leq A t\log t$ for all $t\geq H\geq100$, then
\begin{equation}
 D_{\rm cont}^2\leq
 \frac{\frac{64}{3}A^3(\log H+1)^3
                  +20A(\log H/3+1/9)}{H^3}.
 \label{eq:ns77-height-bound}
\end{equation}
No tail prefix is substituted for the complete zero sums.
\end{proposition}
\begin{proof}
For one pair write $a=\beta$, $c=1-\beta$, $\delta=a-c=2\beta-1$,
so $a+c=1$. Its real-normalized inner factor at $1$ is
$B_\rho(1)=(\gamma^2+c^2)/(\gamma^2+a^2)$.
Set $j_\rho=-\log B_\rho(1)$ and
$\ell_\rho=(\log B_\rho)'(1)$. Exactly,
\[
 j_\rho=\int_{c^2}^{a^2}\frac{dt}{\gamma^2+t},\qquad
 \ell_\rho=-\frac{2\delta(\gamma^2-ac)}
                        {(\gamma^2+a^2)(\gamma^2+c^2)}.
\]
Thus $0\leq j_\rho\leq\delta/\gamma^2$ and
\[
 \ell_\rho+\frac{2\delta}{\gamma^2}
 =\frac{2\delta\{\gamma^2(a^2+c^2+ac)+a^2c^2\}}
        {\gamma^2(\gamma^2+a^2)(\gamma^2+c^2)}
 \leq\frac{17\delta}{8\gamma^4},
\]
using $a^2+c^2+ac\leq1$, $a^2c^2\leq1/16$, and $\gamma\geq1$.
In particular $\ell_\rho+2j_\rho\leq17\delta/(8\gamma^4)$.
Absolute convergence permits summing these inequalities and the
logarithmic derivatives. With
\[
 J=-\log B(1)\in[0,S_2],\qquad \ell=B'(1)/B(1),
 \qquad \epsilon=\ell+2J,
\]
we have $\epsilon\leq17S_4/8$. There is no lower-sign assumption on
$\epsilon$.

Let $\ell_2=(\log B)''(1)$, a real number. Substituting
$b=e^{-J}$, $d=b\ell$, $e=b(\ell^2+\ell_2)$ into
\eqref{eq:ns76-continuous-distance} and discarding only the
nonnegative square that is subtracted gives
\begin{align*}
 D_{\rm cont}^2
 &\leq 2-e^{-2J}(2-2\ell+\ell^2)\\
 &=2-e^{-2J}(2+4J+4J^2)
       +(2+4J)e^{-2J}\epsilon-e^{-2J}\epsilon^2.
\end{align*}
The first term on the last line is zero at $J=0$ and has derivative
$8J^2e^{-2J}$, so it is at most $8J^3/3$.
Also $(2+4J)e^{-2J}\leq2$ for $J\geq0$. If $\epsilon\leq0$ its
linear contribution is nonpositive; otherwise it is at most
$17S_4/4<5S_4$. This proves~\eqref{eq:ns77-tail-budget}.

For the second assertion, $0<2\beta-1<1$ and partial summation give,
for $k=2,4$,
\[
 S_k\leq k\int_H^\infty\frac{N(t)}{t^{k+1}}\,dt
 \leq kA H^{1-k}
             \left(\frac{\log H}{k-1}+\frac1{(k-1)^2}\right).
\]
The endpoint at infinity vanishes and dropping the nonpositive lower
endpoint term only enlarges the bound. Inserting these complete
integrals proves~\eqref{eq:ns77-height-bound}.
\end{proof}

\paragraph{Calibration using published finite-height information.}
The bound of Hasanalizade--Shen--Wong
\cite[Corollary 1.2]{HasanalizadeShenWong2021} implies
$N(t)\leq t\log t$ for $t\geq100$. Indeed its main term is at most
$(t/6)\log t$, and its remainder is at most
$.362\log t+10$; since $\log100>4$, their sum is less than
$(t/6+3)\log t\leq t\log t$ on this range.
Platt--Trudgian's published interval verification
\cite[Theorem 1]{PlattTrudgian2021} supplies the absence of off-critical
zeros up to $H=3\cdot10^{12}$. Taking $A=1$ in the analytic bound
therefore gives the explicit consequence
\begin{equation}
 0\leq D_{\rm cont}^2<2.08\cdot10^{-32}.
 \label{eq:ns77-published-calibration}
\end{equation}
Only the final scalar evaluation is replayed here at 256/384 bits;
the cited large zero verification is an external published input,
not a new computation or an archived verification by this project.
This bound is deliberately conservative.

A small positive continuous defect remains compatible with this
inequality. Moreover the integer and continuous closures are strictly
different by Proposition~\ref{prop:ns76-strict-closure}, and no matching
upper bound on the integer limiting error has been supplied.
Thus the calibration explains why the relaxed problem can look
extremely favorable without proving the original approximation target,
a new zero-free region, G2 or RH.
